% ============================================================
% DOCUMENTACIÓN TÉCNICA — Sistema de Gestión de Mantenimiento
% de Equipos (SGME)
% ============================================================
\documentclass[12pt, a4paper]{article}

% ── Paquetes ─────────────────────────────────────────────────
\usepackage[utf8]{inputenc}
\usepackage[spanish]{babel}
\usepackage[T1]{fontenc}
\usepackage{geometry}
\usepackage{graphicx}
\usepackage{xcolor}
\usepackage{listings}
\usepackage{booktabs}
\usepackage{tabularx}
\usepackage{longtable}
\usepackage{hyperref}
\usepackage{tikz}
\usepackage{float}
\usepackage{fancyhdr}
\usepackage{titlesec}
\usepackage{enumitem}
\usepackage{caption}
\usepackage{tcolorbox}

\usetikzlibrary{shapes.geometric, arrows.meta, positioning, calc, fit}

% ── Configuración de página ──────────────────────────────────
\geometry{margin=2.5cm}

% ── Colores personalizados ───────────────────────────────────
\definecolor{primary}{HTML}{6C5CE7}
\definecolor{secondary}{HTML}{00CEC9}
\definecolor{darkbg}{HTML}{1A1A2E}
\definecolor{codebg}{HTML}{F5F5F5}
\definecolor{codeframe}{HTML}{DDDDDD}
\definecolor{accentgreen}{HTML}{00B894}
\definecolor{accentorange}{HTML}{F39C12}
\definecolor{accentred}{HTML}{E74C3C}

% ── Estilo de código ─────────────────────────────────────────
\lstset{
    backgroundcolor=\color{codebg},
    frame=single,
    rulecolor=\color{codeframe},
    basicstyle=\ttfamily\small,
    keywordstyle=\color{primary}\bfseries,
    commentstyle=\color{accentgreen}\itshape,
    stringstyle=\color{accentorange},
    breaklines=true,
    numbers=left,
    numberstyle=\tiny\color{gray},
    numbersep=8pt,
    tabsize=4,
    showstringspaces=false,
    captionpos=b,
    xleftmargin=1.5em,
    framexleftmargin=1.5em,
    aboveskip=1em,
    belowskip=1em,
}

\lstdefinestyle{sql}{language=SQL, morekeywords={AUTOINCREMENT, PRAGMA, TEXT, REAL, IF}}
\lstdefinestyle{python}{language=Python, morekeywords={self, True, False, None}}
\lstdefinestyle{bash}{language=bash, morekeywords={python, pip, cd}}

% ── Estilo de hipervínculos ──────────────────────────────────
\hypersetup{
    colorlinks=true,
    linkcolor=primary,
    urlcolor=secondary,
    citecolor=primary,
}

% ── Estilo de encabezados ────────────────────────────────────
\pagestyle{fancy}
\fancyhf{}
\fancyhead[L]{\small\textcolor{gray}{SGME — Documentación Técnica}}
\fancyhead[R]{\small\textcolor{gray}{\thepage}}
\fancyfoot[C]{\small\textcolor{gray}{Sistema de Gestión de Mantenimiento de Equipos}}
\renewcommand{\headrulewidth}{0.4pt}

% ── Cajas de información ─────────────────────────────────────
\newtcolorbox{infobox}[1][]{
    colback=primary!5, colframe=primary!60,
    fonttitle=\bfseries, title=#1,
    arc=3mm, boxrule=0.5mm
}

\newtcolorbox{warnbox}[1][]{
    colback=accentorange!5, colframe=accentorange!60,
    fonttitle=\bfseries, title=#1,
    arc=3mm, boxrule=0.5mm
}

% ══════════════════════════════════════════════════════════════
% INICIO DEL DOCUMENTO
% ══════════════════════════════════════════════════════════════
\begin{document}

% ── Portada ──────────────────────────────────────────────────
\begin{titlepage}
    \centering
    \vspace*{2cm}

    {\Huge\bfseries\textcolor{primary}{Sistema de Gestión de}\par}
    \vspace{0.3cm}
    {\Huge\bfseries\textcolor{primary}{Mantenimiento de Equipos}\par}
    \vspace{0.5cm}
    {\Large\textcolor{secondary}{SGME}\par}

    \vspace{0.8cm}
    {\large Documentación Técnica\par}

    \vspace{1.5cm}

    \begin{tikzpicture}
        \draw[primary, line width=2pt] (0,0) -- (12,0);
    \end{tikzpicture}

    \vspace{1.2cm}

    {\large\textbf{Tecnologías:} Python 3.8+ · SQLite · Flask\par}
    \vspace{0.4cm}
    {\large\textbf{Arquitectura:} Modelo--Vista--Controlador (MVC)\par}
    \vspace{0.4cm}
    {\large\textbf{Base de datos:} SQLite 3\par}

    \vspace{1.5cm}

    \begin{tikzpicture}
        \draw[secondary, line width=1pt] (0,0) -- (12,0);
    \end{tikzpicture}

    \vspace{1cm}

    {\large\textbf{Presentado por:}\par}
    \vspace{0.4cm}
    {\Large Emanuel Cardenas\par}
    \vspace{0.2cm}
    {\Large Samuel Manchola\par}
    \vspace{0.2cm}
    {\Large Sebastian Medina\par}

    \vfill
    {\large Mayo 2026\par}
\end{titlepage}

% ── Tabla de Contenidos ──────────────────────────────────────
\tableofcontents
\newpage

% ══════════════════════════════════════════════════════════════
% SECCIÓN 1: INTRODUCCIÓN
% ══════════════════════════════════════════════════════════════
\section{Introducción}

El \textbf{Sistema de Gestión de Mantenimiento de Equipos (SGME)} es una aplicación de software funcional desarrollada en \textbf{Python} con \textbf{SQLite} que gestiona el ciclo de vida completo del mantenimiento de equipos industriales, incluyendo planificación, ejecución, historial y generación de reportes.

\subsection{Objetivos del Sistema}

\begin{itemize}[leftmargin=2em]
    \item Gestionar el inventario de equipos industriales con su ciclo de vida.
    \item Administrar el registro de técnicos especializados.
    \item Controlar las órdenes de trabajo desde su creación hasta su resolución.
    \item Mantener un historial completo de mantenimientos realizados.
    \item Generar reportes especiales para la toma de decisiones.
    \item Aplicar reglas de negocio que garanticen la integridad de los datos.
\end{itemize}

\subsection{Alcance}

El sistema ofrece dos interfaces de usuario:
\begin{enumerate}[leftmargin=2em]
    \item \textbf{Interfaz de consola} (\texttt{main.py}): Menús interactivos basados en texto.
    \item \textbf{Interfaz web} (\texttt{app.py}): Aplicación web moderna con Flask, accesible desde el navegador.
\end{enumerate}

Ambas interfaces comparten la misma base de datos y lógica de negocio.

% ══════════════════════════════════════════════════════════════
% SECCIÓN 2: DIAGRAMA ENTIDAD–RELACIÓN
% ══════════════════════════════════════════════════════════════
\newpage
\section{Diagrama Entidad–Relación}

\subsection{Modelo Conceptual}

El sistema consta de tres entidades principales relacionadas entre sí:

\begin{figure}[H]
    \centering
    \resizebox{\textwidth}{!}{%
    \begin{tikzpicture}[
        entity/.style={
            rectangle, draw=primary, fill=primary!8,
            minimum width=3.8cm, minimum height=1cm,
            font=\bfseries, rounded corners=3pt,
            line width=1pt
        },
        attribute/.style={
            rectangle, draw=gray!50, fill=gray!5,
            minimum width=2.8cm, minimum height=0.5cm,
            font=\small\ttfamily, rounded corners=2pt,
            line width=0.5pt
        },
        pk/.style={
            attribute, draw=accentorange!70, fill=accentorange!8,
            font=\small\ttfamily\bfseries
        },
        relationship/.style={
            diamond, draw=secondary, fill=secondary!10,
            minimum width=1.8cm, minimum height=0.9cm,
            font=\bfseries\small, aspect=2,
            line width=1pt
        },
        line/.style={draw=gray!70, line width=0.8pt, -{Stealth[length=3mm]}},
        cardlabel/.style={font=\footnotesize\bfseries, text=accentred}
    ]

    % ── Entidad EQUIPOS (izquierda) ──
    \node[entity] (equipos) at (0, 0) {EQUIPOS};
    \node[pk]        (eq_pk) at (0, -1.1) {codigo (PK)};
    \node[attribute] (eq_a1) at (0, -1.8) {nombre};
    \node[attribute] (eq_a2) at (0, -2.5) {tipo};
    \node[attribute] (eq_a3) at (0, -3.2) {ubicacion};
    \node[attribute] (eq_a4) at (0, -3.9) {fecha\_instalacion};
    \node[attribute] (eq_a5) at (0, -4.6) {estado};
    \draw[gray!40] (equipos.south) -- (eq_pk.north);
    \draw[gray!40] (eq_pk.south) -- (eq_a1.north);
    \draw[gray!40] (eq_a1.south) -- (eq_a2.north);
    \draw[gray!40] (eq_a2.south) -- (eq_a3.north);
    \draw[gray!40] (eq_a3.south) -- (eq_a4.north);
    \draw[gray!40] (eq_a4.south) -- (eq_a5.north);

    % ── Entidad TÉCNICOS (derecha) ──
    \node[entity] (tecnicos) at (10, 0) {TECNICOS};
    \node[pk]        (tc_pk) at (10, -1.1) {documento (PK)};
    \node[attribute] (tc_a1) at (10, -1.8) {nombre};
    \node[attribute] (tc_a2) at (10, -2.5) {especialidad};
    \node[attribute] (tc_a3) at (10, -3.2) {telefono};
    \node[attribute] (tc_a4) at (10, -3.9) {email (UNIQUE)};
    \draw[gray!40] (tecnicos.south) -- (tc_pk.north);
    \draw[gray!40] (tc_pk.south) -- (tc_a1.north);
    \draw[gray!40] (tc_a1.south) -- (tc_a2.north);
    \draw[gray!40] (tc_a2.south) -- (tc_a3.north);
    \draw[gray!40] (tc_a3.south) -- (tc_a4.north);

    % ── Relaciones (centro) ──
    \node[relationship] (rel1) at (3, -6) {tiene};
    \node[relationship] (rel2) at (7, -6) {asignado};

    % ── Entidad ORDENES DE TRABAJO (centro abajo) ──
    \node[entity, minimum width=5cm] (ordenes) at (5, -8) {ORDENES DE TRABAJO};
    \node[pk]        (ot_pk) at (5, -9.0) {id (PK, AUTO)};
    \node[attribute] (ot_a1) at (3, -9.8) {fecha\_solicitud};
    \node[attribute] (ot_a2) at (3, -10.5) {descripcion\_falla};
    \node[attribute] (ot_a3) at (3, -11.2) {prioridad};
    \node[attribute] (ot_a4) at (3, -11.9) {estado};
    \node[attribute] (ot_a5) at (7, -9.8) {fecha\_ejecucion};
    \node[attribute] (ot_a6) at (7, -10.5) {costo\_repuestos};
    \node[attribute] (ot_a7) at (7, -11.2) {solucion\_aplicada};
    \draw[gray!40] (ordenes.south) -- (ot_pk.north);
    \draw[gray!40] (ot_pk.south) -- ++(0,-0.15) -| (ot_a1.north);
    \draw[gray!40] (ot_a1.south) -- (ot_a2.north);
    \draw[gray!40] (ot_a2.south) -- (ot_a3.north);
    \draw[gray!40] (ot_a3.south) -- (ot_a4.north);
    \draw[gray!40] (ot_pk.south) -- ++(0,-0.15) -| (ot_a5.north);
    \draw[gray!40] (ot_a5.south) -- (ot_a6.north);
    \draw[gray!40] (ot_a6.south) -- (ot_a7.north);

    % ── Lineas de relacion ──
    \draw[line] (equipos.south) -- ++(0,-0.8) -| (rel1.north);
    \draw[line] (rel1.south) -- ++(0,-0.5) -| (ordenes.west);
    \draw[line] (tecnicos.south) -- ++(0,-0.8) -| (rel2.north);
    \draw[line] (rel2.south) -- ++(0,-0.5) -| (ordenes.east);

    % ── Cardinalidad ──
    \node[cardlabel] at (0.5, -5.2) {1};
    \node[cardlabel] at (3, -7.3) {N};
    \node[cardlabel] at (9.5, -5.2) {1};
    \node[cardlabel] at (7, -7.3) {N};

    % ── FK labels ──
    \node[font=\scriptsize\itshape, text=primary] at (1.5, -7.0) {FK: equipo\_codigo};
    \node[font=\scriptsize\itshape, text=primary] at (8.5, -7.0) {FK: tecnico\_documento};

    \end{tikzpicture}
    }% end resizebox
    \caption{Diagrama Entidad--Relación del Sistema SGME}
    \label{fig:er-diagram}
\end{figure}

\subsection{Descripción de las Relaciones}

\begin{table}[H]
    \centering
    \begin{tabularx}{\textwidth}{l l l X}
        \toprule
        \textbf{Relación} & \textbf{Cardinalidad} & \textbf{FK} & \textbf{Restricción} \\
        \midrule
        Equipos $\rightarrow$ Órdenes & 1:N & \texttt{equipo\_codigo} & \texttt{ON DELETE RESTRICT} \\
        Técnicos $\rightarrow$ Órdenes & 1:N & \texttt{tecnico\_documento} & \texttt{ON DELETE RESTRICT} \\
        \bottomrule
    \end{tabularx}
    \caption{Relaciones entre entidades}
\end{table}

\begin{itemize}[leftmargin=2em]
    \item Un \textbf{equipo} puede tener muchas órdenes de trabajo (1:N).
    \item Un \textbf{técnico} puede ser asignado a muchas órdenes de trabajo (1:N).
    \item Una \textbf{orden de trabajo} pertenece a exactamente un equipo y un técnico.
    \item Las claves foráneas usan \texttt{ON DELETE RESTRICT} para impedir la eliminación de equipos o técnicos con órdenes asociadas.
\end{itemize}

\subsection{Restricciones de Integridad}

\begin{table}[H]
    \centering
    \begin{tabularx}{\textwidth}{l l X}
        \toprule
        \textbf{Tabla} & \textbf{Tipo} & \textbf{Restricción} \\
        \midrule
        \texttt{equipos} & CHECK & \texttt{estado IN ('operativo', 'en mantenimiento', 'fuera de servicio')} \\
        \texttt{ordenes\_trabajo} & CHECK & \texttt{prioridad IN ('alta', 'media', 'baja')} \\
        \texttt{ordenes\_trabajo} & CHECK & \texttt{estado IN ('pendiente', 'en proceso', 'completada', 'cancelada')} \\
        \texttt{tecnicos} & UNIQUE & \texttt{email} debe ser único \\
        \bottomrule
    \end{tabularx}
    \caption{Restricciones de integridad referencial y de dominio}
\end{table}

\subsection{Índices}

\begin{table}[H]
    \centering
    \begin{tabularx}{\textwidth}{l X l}
        \toprule
        \textbf{Índice} & \textbf{Columna(s)} & \textbf{Propósito} \\
        \midrule
        \texttt{idx\_equipos\_estado} & \texttt{equipos(estado)} & Filtrar equipos por estado \\
        \texttt{idx\_ordenes\_equipo} & \texttt{ordenes\_trabajo(equipo\_codigo)} & JOINs con equipos \\
        \texttt{idx\_ordenes\_tecnico} & \texttt{ordenes\_trabajo(tecnico\_documento)} & JOINs con técnicos \\
        \texttt{idx\_ordenes\_estado} & \texttt{ordenes\_trabajo(estado)} & Filtrar OTs por estado \\
        \texttt{idx\_ordenes\_fecha} & \texttt{ordenes\_trabajo(fecha\_solicitud)} & Consultas por fecha \\
        \bottomrule
    \end{tabularx}
    \caption{Índices definidos para optimización de consultas}
\end{table}

% ══════════════════════════════════════════════════════════════
% SECCIÓN 3: ARQUITECTURA DEL SISTEMA
% ══════════════════════════════════════════════════════════════
\newpage
\section{Arquitectura del Sistema}

El sistema implementa una arquitectura de \textbf{tres capas} basada en el patrón \textbf{Modelo--Vista--Controlador (MVC)}, separando claramente las responsabilidades de datos, lógica de negocio y presentación.

\subsection{Diagrama de Capas}

\begin{figure}[H]
    \centering
    \begin{tikzpicture}[
        layer/.style={
            rectangle, rounded corners=5pt, minimum width=10cm,
            minimum height=1.2cm, line width=1pt, font=\bfseries,
            text=white, align=center
        },
        module/.style={
            rectangle, rounded corners=3pt,
            minimum height=0.6cm, font=\footnotesize\ttfamily, line width=0.5pt,
            inner sep=4pt
        },
        dbnode/.style={
            rectangle, rounded corners=5pt, minimum width=3.5cm,
            minimum height=0.8cm, line width=1pt, font=\bfseries\footnotesize,
            draw=accentred, fill=accentred!10, text=accentred
        },
        arrow/.style={
            -{Stealth[length=2.5mm]}, line width=1pt, gray!60
        }
    ]

    % ── CAPA 1: Presentacion ──
    \node[layer, fill=primary] at (0, 0) {Capa de Presentacion};
    \node[module, fill=primary!15, draw=primary!40, text=darkbg]
        at (-3, -1.0) {views/menu.py};
    \node[module, fill=primary!15, draw=primary!40, text=darkbg]
        at (0, -1.0) {views/tabla.py};
    \node[module, fill=primary!15, draw=primary!40, text=darkbg]
        at (3, -1.0) {templates/*.html};
    \node[module, fill=primary!25, draw=primary!50, text=darkbg]
        at (-3.5, -1.8) {app.py {\tiny(Flask)}};
    \node[module, fill=primary!25, draw=primary!50, text=darkbg]
        at (3.5, -1.8) {main.py {\tiny(Consola)}};

    \draw[arrow] (0, -2.2) -- (0, -2.8);

    % ── CAPA 2: Logica de Negocio ──
    \node[layer, fill=secondary!80!black] at (0, -3.5) {Capa de Logica de Negocio};
    \node[module, fill=secondary!10, draw=secondary!40, text=darkbg]
        at (-3.5, -4.5) {models/equipo.py};
    \node[module, fill=secondary!10, draw=secondary!40, text=darkbg]
        at (0, -4.5) {models/tecnico.py};
    \node[module, fill=secondary!10, draw=secondary!40, text=darkbg]
        at (3.7, -4.5) {models/orden\_trabajo.py};
    \node[module, fill=accentgreen!10, draw=accentgreen!40, text=darkbg]
        at (-2, -5.3) {reports/historial.py};
    \node[module, fill=accentgreen!10, draw=accentgreen!40, text=darkbg]
        at (2, -5.3) {reports/reportes.py};

    \draw[arrow] (0, -5.8) -- (0, -6.4);

    % ── CAPA 3: Datos ──
    \node[layer, fill=accentorange!80!black] at (0, -7.1) {Capa de Datos};
    \node[module, fill=accentorange!10, draw=accentorange!40, text=darkbg]
        at (-2, -8.1) {database.py};
    \node[module, fill=accentorange!10, draw=accentorange!40, text=darkbg]
        at (2, -8.1) {schema.sql};

    \draw[arrow] (0, -8.5) -- (0, -9.1);

    % ── Base de Datos ──
    \node[dbnode] at (0, -9.7) {mantenimiento.db};

    \end{tikzpicture}
    \caption{Arquitectura de tres capas del sistema SGME}
    \label{fig:architecture}
\end{figure}

\subsection{Descripción de las Capas}

\subsubsection{Capa de Presentación (Vista)}

Es la capa que interactúa directamente con el usuario. Existen dos implementaciones:

\begin{itemize}[leftmargin=2em]
    \item \textbf{Interfaz Web (Flask)}: Plantillas HTML con CSS moderno, accesible desde \texttt{http://localhost:5000}. Los archivos se encuentran en \texttt{templates/} y \texttt{static/}.
    \item \textbf{Interfaz de Consola}: Menús interactivos con tablas ASCII formateadas, implementados en \texttt{views/menu.py} y \texttt{views/tabla.py}.
\end{itemize}

Esta capa \textbf{no contiene lógica de negocio}; solo se encarga de mostrar información y capturar la entrada del usuario.

\subsubsection{Capa de Lógica de Negocio (Modelo)}

Contiene todas las reglas de negocio y operaciones CRUD:

\begin{itemize}[leftmargin=2em]
    \item \textbf{Modelos} (\texttt{models/}): Operaciones de creación, lectura, actualización y eliminación de equipos, técnicos y órdenes de trabajo. Aquí se validan las reglas de negocio.
    \item \textbf{Reportes} (\texttt{reports/}): Consultas SQL especializadas para historial y reportes avanzados.
\end{itemize}

\subsubsection{Capa de Datos}

Gestiona la persistencia y el acceso a la base de datos:

\begin{itemize}[leftmargin=2em]
    \item \textbf{\texttt{database.py}}: Clase Singleton que gestiona la conexión SQLite, ejecuta consultas y controla transacciones.
    \item \textbf{\texttt{schema.sql}}: Script DDL con la definición de tablas, índices, restricciones y datos de ejemplo.
    \item \textbf{\texttt{mantenimiento.db}}: Archivo de base de datos SQLite generado automáticamente.
\end{itemize}

% ══════════════════════════════════════════════════════════════
% SECCIÓN 4: REGLAS DE NEGOCIO
% ══════════════════════════════════════════════════════════════
\subsection{Reglas de Negocio}

\begin{table}[H]
    \centering
    \begin{tabularx}{\textwidth}{l l X}
        \toprule
        \textbf{ID} & \textbf{Módulo} & \textbf{Descripción} \\
        \midrule
        RN-01 & \texttt{equipo.py} & No se puede eliminar un equipo que tenga órdenes de trabajo asociadas. Se verifica mediante \texttt{COUNT} antes del \texttt{DELETE}, reforzado por \texttt{FK ON DELETE RESTRICT}. \\
        \addlinespace
        RN-02 & \texttt{orden\_trabajo.py} & No se puede asignar un técnico que no esté registrado en el sistema. Se verifica la existencia del técnico en la tabla antes de crear la orden. \\
        \addlinespace
        RN-03 & \texttt{orden\_trabajo.py} & Al completar una orden de trabajo, el equipo asociado cambia automáticamente a estado ``operativo''. Se ejecuta dentro de la misma transacción. \\
        \bottomrule
    \end{tabularx}
    \caption{Reglas de negocio implementadas}
\end{table}

% ══════════════════════════════════════════════════════════════
% SECCIÓN 5: CONFIGURACIÓN E INSTALACIÓN
% ══════════════════════════════════════════════════════════════
\newpage
\section{Configuración Requerida}

\subsection{Requisitos del Sistema}

\begin{table}[H]
    \centering
    \begin{tabularx}{\textwidth}{l X}
        \toprule
        \textbf{Componente} & \textbf{Requisito} \\
        \midrule
        Sistema Operativo & Windows 10/11, Linux o macOS \\
        Python & Versión \textbf{3.8} o superior \\
        SQLite & Incluido con Python (módulo \texttt{sqlite3}) \\
        Flask & Versión 3.x (para la interfaz web) \\
        Navegador & Cualquier navegador moderno (Chrome, Firefox, Edge) \\
        \bottomrule
    \end{tabularx}
    \caption{Requisitos mínimos del sistema}
\end{table}

\subsection{Dependencias}

\begin{table}[H]
    \centering
    \begin{tabularx}{\textwidth}{l l X}
        \toprule
        \textbf{Paquete} & \textbf{Tipo} & \textbf{Uso} \\
        \midrule
        \texttt{sqlite3} & Estándar (incluido) & Motor de base de datos \\
        \texttt{os} & Estándar (incluido) & Operaciones del sistema de archivos \\
        \texttt{sys} & Estándar (incluido) & Configuración del intérprete \\
        \texttt{datetime} & Estándar (incluido) & Manejo de fechas \\
        \texttt{flask} & Externo (\texttt{pip}) & Framework web para la interfaz gráfica \\
        \bottomrule
    \end{tabularx}
    \caption{Dependencias del proyecto}
\end{table}

\begin{infobox}[Nota]
Para la interfaz de consola (\texttt{main.py}), \textbf{no se requiere ninguna dependencia externa}. Todo funciona con la librería estándar de Python.
\end{infobox}

\subsection{Instrucciones de Instalación}

\begin{enumerate}[leftmargin=2em]
    \item \textbf{Verificar versión de Python:}
    \begin{lstlisting}[style=bash]
python --version
# Debe mostrar Python 3.8 o superior
    \end{lstlisting}

    \item \textbf{Instalar Flask} (solo para interfaz web):
    \begin{lstlisting}[style=bash]
pip install flask
    \end{lstlisting}

    \item \textbf{Ubicarse en el directorio del proyecto:}
    \begin{lstlisting}[style=bash]
cd "proyecto final paradigmas"
    \end{lstlisting}
\end{enumerate}

\subsection{Instrucciones de Ejecución}

\subsubsection{Opción 1: Interfaz Web (recomendada)}

\begin{lstlisting}[style=bash]
python app.py
\end{lstlisting}

Al ejecutar, el sistema:
\begin{enumerate}[leftmargin=2em]
    \item Inicializa la base de datos automáticamente (si no existe).
    \item Inicia el servidor web en \texttt{http://localhost:5000}.
    \item Abrir esa URL en el navegador para acceder al sistema.
\end{enumerate}

\subsubsection{Opción 2: Interfaz de Consola}

\begin{lstlisting}[style=bash]
python main.py
\end{lstlisting}

Muestra un menú interactivo en la terminal con opciones numeradas.

\begin{warnbox}[Importante]
La primera ejecución crea el archivo \texttt{mantenimiento.db} y carga los datos de ejemplo automáticamente. Las ejecuciones posteriores detectan que la base de datos ya existe y no duplican los datos.
\end{warnbox}

% ══════════════════════════════════════════════════════════════
% SECCIÓN 6: DESCRIPCIÓN DE MÓDULOS
% ══════════════════════════════════════════════════════════════
\newpage
\section{Descripción de Módulos}

\subsection{Estructura de Archivos}

\begin{lstlisting}[language={}, basicstyle=\ttfamily\small, frame=none, numbers=none]
proyecto final paradigmas/
|-- schema.sql                 # Script DDL + datos de ejemplo
|-- database.py                # Conexion y gestion de BD (Singleton)
|-- main.py                    # Punto de entrada (consola)
|-- app.py                     # Punto de entrada (web Flask)
|-- models/
|   |-- __init__.py
|   |-- equipo.py              # CRUD de equipos
|   |-- tecnico.py             # CRUD de tecnicos
|   |-- orden_trabajo.py       # CRUD avanzado de ordenes
|-- views/
|   |-- __init__.py
|   |-- menu.py                # Menus interactivos de consola
|   |-- tabla.py               # Formateo de tablas ASCII
|-- reports/
|   |-- __init__.py
|   |-- historial.py           # Historial de mantenimiento
|   |-- reportes.py            # Reportes especiales
|-- templates/                 # Plantillas HTML (Flask)
|   |-- base.html              # Layout base con sidebar
|   |-- dashboard.html         # Pagina principal
|   |-- equipos/               # CRUD de equipos
|   |-- tecnicos/              # CRUD de tecnicos
|   |-- ordenes/               # CRUD de ordenes
|   |-- historial.html         # Historial con filtros
|   |-- reportes.html          # Reportes especiales
|-- static/
|   |-- style.css              # Estilos CSS (diseno oscuro)
|-- mantenimiento.db           # Base de datos SQLite (autogenerada)
\end{lstlisting}

\subsection{Módulo: \texttt{schema.sql}}

\begin{table}[H]
    \centering
    \begin{tabularx}{\textwidth}{l X}
        \toprule
        \textbf{Responsabilidad} & Definición del esquema de base de datos \\
        \textbf{Tipo} & Script SQL \\
        \midrule
        \textbf{Contenido} & Sentencias \texttt{CREATE TABLE} para las 3 tablas del sistema, \texttt{CREATE INDEX} para los 5 índices de optimización, restricciones \texttt{CHECK} y \texttt{FOREIGN KEY}, e instrucciones \texttt{INSERT} con datos de ejemplo (15 equipos, 12 técnicos, 15 órdenes). \\
        \bottomrule
    \end{tabularx}
\end{table}

\subsection{Módulo: \texttt{database.py}}

\begin{table}[H]
    \centering
    \begin{tabularx}{\textwidth}{l X}
        \toprule
        \textbf{Responsabilidad} & Gestión de la conexión y operaciones sobre la base de datos \\
        \textbf{Patrón} & Singleton \\
        \midrule
        \textbf{Clase} & \texttt{Database} \\
        \textbf{Métodos} & \texttt{conectar()}: Establece conexión SQLite con \texttt{PRAGMA foreign\_keys = ON} y \texttt{check\_same\_thread=False}. \\
        & \texttt{inicializar()}: Ejecuta \texttt{schema.sql} si las tablas no existen. \\
        & \texttt{ejecutar(query, params)}: INSERT/UPDATE/DELETE con commit y rollback. \\
        & \texttt{fetchone(query, params)}: Retorna un registro. \\
        & \texttt{fetchall(query, params)}: Retorna todos los registros. \\
        & \texttt{cerrar()}: Cierra la conexión. \\
        \bottomrule
    \end{tabularx}
\end{table}

\subsection{Módulo: \texttt{models/equipo.py}}

\begin{table}[H]
    \centering
    \begin{tabularx}{\textwidth}{l X}
        \toprule
        \textbf{Responsabilidad} & Operaciones CRUD para la entidad Equipos \\
        \textbf{Regla de negocio} & RN-01: No eliminar equipos con OTs asociadas \\
        \midrule
        \textbf{Funciones} & \texttt{crear\_equipo(codigo, nombre, tipo, ubicacion, fecha, estado)}: Crea un equipo validando el estado. \\
        & \texttt{obtener\_equipos()}: Retorna todos los equipos ordenados por código. \\
        & \texttt{obtener\_equipo(codigo)}: Busca un equipo por su clave primaria. \\
        & \texttt{actualizar\_equipo(codigo, **campos)}: Actualiza campos específicos. \\
        & \texttt{eliminar\_equipo(codigo)}: Verifica OTs asociadas antes de eliminar. \\
        \bottomrule
    \end{tabularx}
\end{table}

\subsection{Módulo: \texttt{models/tecnico.py}}

\begin{table}[H]
    \centering
    \begin{tabularx}{\textwidth}{l X}
        \toprule
        \textbf{Responsabilidad} & Operaciones CRUD para la entidad Técnicos \\
        \textbf{Validaciones} & Email único, protección contra eliminación con OTs activas \\
        \midrule
        \textbf{Funciones} & \texttt{crear\_tecnico(documento, nombre, especialidad, tel, email)}: Crea un técnico con validación de unicidad. \\
        & \texttt{obtener\_tecnicos()}: Lista todos los técnicos. \\
        & \texttt{obtener\_tecnico(documento)}: Busca por documento. \\
        & \texttt{actualizar\_tecnico(documento, **campos)}: Actualiza datos. \\
        & \texttt{eliminar\_tecnico(documento)}: Verifica OTs activas antes de eliminar. \\
        \bottomrule
    \end{tabularx}
\end{table}

\subsection{Módulo: \texttt{models/orden\_trabajo.py}}

\begin{table}[H]
    \centering
    \begin{tabularx}{\textwidth}{l X}
        \toprule
        \textbf{Responsabilidad} & CRUD avanzado de Órdenes de Trabajo con reglas de negocio \\
        \textbf{Reglas} & RN-02 (técnico registrado) y RN-03 (equipo operativo al completar) \\
        \midrule
        \textbf{Funciones} & \texttt{crear\_orden(...)}: Valida equipo y técnico, crea OT, cambia equipo a ``en mantenimiento''. \\
        & \texttt{obtener\_ordenes()}: Lista OTs con JOIN a equipos y técnicos. \\
        & \texttt{obtener\_orden(id)}: Detalle completo de una OT. \\
        & \texttt{actualizar\_orden(id, **campos)}: Modifica campos permitidos. \\
        & \texttt{completar\_orden(id, solucion, costo)}: Cierra la OT, registra fecha, cambia equipo a ``operativo''. \\
        & \texttt{cancelar\_orden(id)}: Cancela una OT (no completadas). \\
        & \texttt{eliminar\_orden(id)}: Solo permite eliminar OTs pendientes o canceladas. \\
        \bottomrule
    \end{tabularx}
\end{table}

\subsection{Módulo: \texttt{reports/historial.py}}

\begin{table}[H]
    \centering
    \begin{tabularx}{\textwidth}{l X}
        \toprule
        \textbf{Responsabilidad} & Consultas de solo lectura sobre el historial de mantenimiento \\
        \midrule
        \textbf{Funciones} & \texttt{ver\_historial(filtro\_equipo, filtro\_tecnico, fecha\_inicio, fecha\_fin)}: Retorna OTs completadas con JOINs a equipos y técnicos, aplicando filtros dinámicos opcionales. Ordenadas por fecha de ejecución descendente. \\
        \bottomrule
    \end{tabularx}
\end{table}

\subsection{Módulo: \texttt{reports/reportes.py}}

\begin{table}[H]
    \centering
    \begin{tabularx}{\textwidth}{l X}
        \toprule
        \textbf{Responsabilidad} & Reportes y consultas especiales avanzadas \\
        \midrule
        \textbf{Funciones} & \texttt{equipos\_sin\_mantenimiento()}: \texttt{LEFT JOIN} para encontrar equipos sin OTs completadas. \\
        & \texttt{top\_equipos\_intervenidos(n=2)}: \texttt{GROUP BY} + \texttt{COUNT} + \texttt{ORDER BY DESC LIMIT 2}. \\
        & \texttt{tecnico\_mayor\_costo()}: \texttt{SUM(costo\_repuestos)} agrupado por técnico. \\
        & \texttt{ordenes\_atrasadas()}: Filtra OTs pendientes/en proceso con más de 7 días usando \texttt{julianday()}. \\
        \bottomrule
    \end{tabularx}
\end{table}

\subsection{Módulo: \texttt{app.py} (Controlador Web)}

\begin{table}[H]
    \centering
    \begin{tabularx}{\textwidth}{l X}
        \toprule
        \textbf{Responsabilidad} & Servidor web Flask — Punto de entrada de la interfaz gráfica \\
        \textbf{Framework} & Flask 3.x \\
        \midrule
        \textbf{Rutas} & \texttt{/} — Dashboard con estadísticas y resumen. \\
        & \texttt{/equipos} — CRUD completo de equipos (listar, crear, editar, eliminar). \\
        & \texttt{/tecnicos} — CRUD completo de técnicos. \\
        & \texttt{/ordenes} — CRUD avanzado de OTs (crear, editar, completar, cancelar, eliminar). \\
        & \texttt{/historial} — Historial con filtros dinámicos por equipo, técnico y fechas. \\
        & \texttt{/reportes} — Los 4 reportes especiales. \\
        \bottomrule
    \end{tabularx}
\end{table}

\subsection{Módulo: \texttt{views/menu.py} (Controlador Consola)}

\begin{table}[H]
    \centering
    \begin{tabularx}{\textwidth}{l X}
        \toprule
        \textbf{Responsabilidad} & Menús interactivos de la interfaz de consola \\
        \midrule
        \textbf{Funciones} & \texttt{menu\_principal()}: Navegación principal con 5 secciones. \\
        & \texttt{menu\_equipos()}: Submenú CRUD de equipos con formularios de entrada. \\
        & \texttt{menu\_tecnicos()}: Submenú CRUD de técnicos. \\
        & \texttt{menu\_ordenes()}: Submenú CRUD avanzado de OTs (8 opciones). \\
        & \texttt{menu\_historial()}: Submenú de historial con 4 tipos de filtro. \\
        & \texttt{menu\_reportes()}: Submenú de reportes especiales. \\
        \bottomrule
    \end{tabularx}
\end{table}

\subsection{Módulo: \texttt{views/tabla.py}}

\begin{table}[H]
    \centering
    \begin{tabularx}{\textwidth}{l X}
        \toprule
        \textbf{Responsabilidad} & Utilidades de formateo y presentación para la consola \\
        \midrule
        \textbf{Funciones} & \texttt{imprimir\_tabla(headers, rows)}: Formatea datos en tabla ASCII alineada. \\
        & \texttt{imprimir\_detalle(campos)}: Muestra un registro en formato clave--valor. \\
        & \texttt{imprimir\_titulo(titulo)}: Encabezado decorado para secciones. \\
        & \texttt{limpiar\_pantalla()}: Limpia la consola (compatible Windows/Linux). \\
        & \texttt{pausar()}: Espera confirmación del usuario (Enter). \\
        & \texttt{confirmar(mensaje)}: Solicita confirmación sí/no. \\
        \bottomrule
    \end{tabularx}
\end{table}

% ══════════════════════════════════════════════════════════════
% SECCIÓN 7: DATOS DE EJEMPLO
% ══════════════════════════════════════════════════════════════
\newpage
\section{Datos de Ejemplo}

El archivo \texttt{schema.sql} incluye datos precargados para pruebas inmediatas:

\begin{table}[H]
    \centering
    \begin{tabularx}{\textwidth}{l c X}
        \toprule
        \textbf{Tabla} & \textbf{Registros} & \textbf{Descripción} \\
        \midrule
        \texttt{equipos} & 15 & Equipos industriales variados (compresores, tornos CNC, bombas, generadores, soldadoras, etc.) con tres estados diferentes. \\
        \addlinespace
        \texttt{tecnicos} & 12 & Técnicos con especialidades diversas (mecánica, electricidad, soldadura, HVAC, automatización, etc.) \\
        \addlinespace
        \texttt{ordenes\_trabajo} & 15 & 9 completadas, 2 en proceso, 3 pendientes (2 atrasadas) y 1 cancelada. Costos variados para reportes significativos. \\
        \bottomrule
    \end{tabularx}
    \caption{Resumen de datos de ejemplo precargados}
\end{table}

Los datos están diseñados para que todos los reportes especiales arrojen resultados significativos:
\begin{itemize}[leftmargin=2em]
    \item Los equipos \texttt{EQ-013}, \texttt{EQ-014} y \texttt{EQ-015} no tienen OTs (reporte de equipos sin mantenimiento).
    \item Los equipos \texttt{EQ-001} y \texttt{EQ-002} tienen múltiples intervenciones (reporte top 2).
    \item Hay OTs con más de 7 días sin completar (reporte de atrasadas).
\end{itemize}

\end{document}
